\documentclass[10pt,twocolumn,letterpaper]{article}
\usepackage[top=2cm,bottom=2.5cm,left=1.8cm,right=1.8cm,columnsep=0.8cm]{geometry}
\usepackage{times}
\usepackage{epsfig}
\usepackage{graphicx}
\usepackage{amsmath}
\usepackage{amssymb}
\usepackage{booktabs}
\usepackage{microtype}
\usepackage{hyperref}
\usepackage{xcolor}
\usepackage{caption}
\usepackage{subcaption}
\hypersetup{colorlinks=true,linkcolor=blue,citecolor=blue,urlcolor=blue}
\title{\textbf{Manual VLA: Methodology and AI Architecture Overview}}
\author{Autonomous Robotics \\& VLA Research Lab\texttt{https://github.com/krazykarthik2/manual_vla}}
\date{September 2026}
\begin{document}
\maketitle
\begin{abstract}
This paper provides a high‑level methodology overview of the repository, summarising the design principles, comparative analysis, and overarching architecture that underpins both the SmolVLA and Full SmolVLA systems.
\end{abstract}
\section{Introduction}
Embodied AI for robotic manipulation demands efficient perception‑action pipelines that can operate in real time. We present two compact multimodal policies—SmolVLA and Full SmolVLA—each built on a distinct backbone but sharing a common Optimal Transport Conditional Flow Matching (OT‑CFM) trajectory generation framework.
\section{System Overview}
Both systems ingest RGB images and natural‑language instructions, encode them via a vision‑language backbone, and generate 128‑step continuous trajectories for a 4‑DOF Dobot Magician robot. Figure~\ref{fig:arch_comparison} (shared from the original manuscript) illustrates the architectural differences.
\begin{figure*}[t]
\centering
\includegraphics[width=\textwidth]{../paper/figures/fig1_architectures.png}
\caption{System Architecture Comparison. Left: SmolVLA dual‑stream pipeline with frozen CLIP ViT‑B/32. Right: Full SmolVLA using the SmolVLM‑256M foundation model.}
\label{fig:arch_comparison}
\end{figure*}
\section{Comparative Analysis}
The two approaches trade off inference latency against perceptual richness. SmolVLA achieves $>100$ Hz inference with a 1.4 M‑parameter decoder, while Full SmolVLA attains higher distractor robustness at $\approx30$ Hz with a 2.1 M‑parameter decoder.
\section{Conclusion}
We demonstrate that compact multimodal architectures can meet real‑time robotic demands without sacrificing performance, providing a practical foundation for edge‑deployed VLA systems.
\bibliographystyle{plain}
\begin{thebibliography}{1}
\bibitem{lipman2022flow}
Yaron Lipman, Ricky T.~Q. Chen, Heli Ben‑Hamu, Maximilian Nickel, and Matt Le.
\newblock Flow matching for generative modeling.
\newblock In {\em ICLR}, 2023.
\bibitem{radford2021learning}
Alec Radford et~al.
\newblock Learning transferable visual models from natural language supervision.
\newblock In {\em ICML}, 2021.
\bib{smolvlm2024}
Hugging Face SmolModels Team.
\newblock SmolVLM: Small, fast, and capable multimodal vision‑language models.
\newblock {\em Hugging Face Technical Report}, 2024.
\end{thebibliography}
\end{document}
